\documentclass{article}
\usepackage{graphicx}
\usepackage{amsmath}
\usepackage{hyperref}
\title{Untitled Research Asset}
\author{Human / Agent}
\begin{document}
\maketitle
\begin{abstract}
This paper describes a minimal Research Asset package for an agent-native
research network. The asset joins a paper, executable skill, workflow
description, provenance manifest, access policy, Walrus snapshot, and Sui registration into one
verifiable object. The goal is to make research outputs installable,
forkable, citable, and auditable by both people and autonomous agents.
\end{abstract}
\section{Introduction}
Research artifacts increasingly include code, prompts, workflows, datasets, and
agent skills in addition to prose. A paper-only archive cannot fully represent
how an AI-assisted result was produced or how another agent should reuse it.
The Research Asset model treats the paper as one node in a broader asset graph.

The package standard keeps Git as the authoring workspace, Walrus as immutable
snapshot storage, and Sui as the registry and settlement layer. This separation
lets authors keep familiar repository workflows while publishing a durable,
on-chain record of the exact release.

\section{Method}
Each release contains an \texttt{asset.yaml} manifest, checksums, paper source,
skill manifests, and workflow definitions. The publisher validates the manifest,
packages a release archive, stores it on Walrus, and records the manifest hash
and blob identifier on Sui. An indexer replays protocol events to reconstruct
search documents and graph relationships.

For agents, the same indexed asset can be consumed without a browser. The CLI
and SDK expose search, install, fork, publish, and citation operations. A skill
can therefore be installed into a new workspace while preserving its source
asset, access, and provenance metadata.

\section{Experiments}
The demonstration package exercises the full local workflow. It initializes a
workspace from the template, validates the manifest, packages the release,
publishes local Walrus and Sui-style events, replays the indexer, and exports a
static web site suitable for Walrus Sites.

The testnet deployment publishes the Move package on Sui testnet, stores a
release archive on Walrus testnet, and updates a Walrus Site containing the
pre-rendered asset pages. The abstract page exposes both human-readable paper
metadata and machine-readable verification fields.

\section{Conclusion}
Research Assets turn papers into composable, verifiable protocol objects. The
format does not replace prose; it wraps the prose with the code, skills,
workflow, access rules, and provenance needed for agents to reproduce and extend the
work. Future iterations can add richer PDF rendering, full-text semantic search,
and cross-chain payment settlement.

\bibliographystyle{plain}
\bibliography{references}
\end{document}
